\documentclass[12pt,letterpaper]{article}

\usepackage[spanish,es-tabla,es-nodecimaldot]{babel}
\usepackage{amsmath}
\usepackage[utf8]{inputenc}
\usepackage[T1]{fontenc}
\usepackage{lmodern}
\usepackage{graphicx}
\usepackage{listings}
\usepackage{anysize} 
\usepackage{fancyhdr}
\usepackage{amsmath}
\usepackage{pdfpages}
\usepackage{graphics}
\usepackage{capt-of}
\usepackage{tabularx}
\usepackage[colorlinks=true,plainpages=true,citecolor=blue,linkcolor=blue]{hyperref}

\marginsize{2cm}{2cm}{2cm}{2cm}
\pagestyle{fancy}
\fancyhf{Líneas de transmisión y antenas}
\fancyhead[L]{\footnotesize UPIITA-IPN} 
\fancyhead[R]{\footnotesize 3TV1} 
\fancyfoot[R]{\footnotesize Práctica 1}
\fancyfoot[C]{\thepage}
\fancyfoot[L]{\footnotesize Carta de Smith} 

\renewcommand{\footrulewidth}{0.4pt}
\renewcommand{\spanishtablename}{Tabla}
\renewcommand{\labelitemii}{$\star$}

\begin{document}

\includepdf[pages={1}]{portada}

%\newpage
%\tableofcontents
%\listoffigures
%\listoftables

\newpage
\section{Balun}
Un balun es un dispositivo que une una línea balanceada (una que tiene dos conductores, 
con corrientes iguales en direcciones opuestas, como un cable de par trenzado) a una 
línea no balanceada (una que tiene sólo un conductor y una tierra, como un cable coaxial). 
Un balun es un tipo de transformador: se utiliza para convertir una señal no balanceada 
en balanceada o viceversa. Los baluns aíslan una línea de transmisión y proporcionan una 
salida balanceada. Un uso típico para un balun es en una antena de televisión. El término 
se deriva de la combinación de equilibrado y desequilibrado.
\\ \\
En un balun, un par de terminales está equilibrado, es decir, las corrientes son iguales 
en magnitud y opuestas en fase. El otro par de terminales está desequilibrado; un lado 
está conectado a tierra eléctrica y el otro lleva la señal.
\\ \\
Hay dos variaciones de este dispositivo:
\begin{itemize}
    \item El unun, que transfiere la señal de una línea desequilibrada a otra.
    \item El balbal, que transfiere la señal de una línea balanceada a otra.
\end{itemize}
Para funcionar con una eficiencia óptima, un balun debe utilizarse con cargas cuyas 
impedancias presentan poca o ninguna reactancia. Tales impedancias son llamadas 
"puramente resistivas".  Como regla general, las antenas de comunicaciones bien 
diseñadas presentan cargas puramente resistivas de 50, 75 o 300 $\Omega$, aunque unas 
pocas antenas tienen impedancias resistivas más altas.
\\ \\
Algunos baluns pueden funcionar como un transformador de impedancia entre dos sistemas 
desequilibrados si hay poca o ninguna reactancia. 

\section{Aplicaciones}
Los transformadores de balun se pueden utilizar entre varias partes de un sistema de 
comunicaciones por cable o inalámbrico. La siguiente tabla denota algunas aplicaciones 
comunes.

\begin{table}[ht]
    \centering
    \begin{tabular}{|l|l|}
    \hline
    \textbf{Equilibrado} & \textbf{Desequilibrado} \\ \hline
    Receptor de televisión & Red de cable coaxial \\ \hline
    Receptor de televisión & Sistema de antena coaxial \\ \hline
    Receptor de radiodifusión FM & Sistema de antena coaxial \\ \hline
    Antena dipolo & Línea de transmisión coaxial \\ \hline
    Línea de transmisión de cable paralelo & Salida del transmisor coaxial \\ \hline
    Línea de transmisión de cable paralelo & Entrada para receptor coaxial \\ \hline
    Línea de transmisión de cable paralelo & Línea de transmisión coaxial \\ \hline
    \end{tabular}
    \caption{Aplicaciones.}
\end{table}

\section{Cálculos}
Impedancia característica de la línea.
\begin{equation}
    Z_0=200 \ \Omega
\end{equation}
Coeficiente de reflexión.
\begin{equation}
    \Gamma(0)=\frac{Z_L-Z_0}{Z_0+Z_L}
\end{equation}
Potencia promedio.
\begin{equation}
    P_{prom}=||\Gamma(0)||^{2}
\end{equation}
Perdida de retorno.
\begin{equation}
    RL=-10\log_{10} (|\Gamma(0)|^{2})
\end{equation}

\subsection{$Z_L=47 \Omega$}
\begin{equation}
    \hat{Z}=\frac{Z_L}{Z_0}=\frac{47}{200}=2.35
\end{equation}
\begin{equation}
    \Gamma(0)=\frac{47-200}{47+200}=-0.6194
\end{equation}
\begin{equation}
    |\Gamma(0)|^{2}=0.836
\end{equation}
\begin{equation}
    RL=-10\log_{10} (-0.6194^{2})=4.1605
\end{equation}

\subsection{$Z_L=220 \Omega$}
\begin{equation}
    \hat{Z}=\frac{Z_L}{Z_0}=\frac{220}{200}=2.35
\end{equation}
\begin{equation}
    \Gamma(0)=\frac{220-200}{220+200}=0.0476
\end{equation}
\begin{equation}
    |\Gamma(0)|^{2}=0.0022
\end{equation}
\begin{equation}
    RL=-10\log_{10} (0.0476^{2})=26.4478
\end{equation}

\subsection{$Z_L=270 \Omega$}
\begin{equation}
    \hat{Z}=\frac{Z_L}{Z_0}=\frac{270}{200}=1.35
\end{equation}
\begin{equation}
    \Gamma(0)=\frac{270-200}{270+200}=0.1489
\end{equation}
\begin{equation}
    |\Gamma(0)|^{2}=0.0221
\end{equation}
\begin{equation}
    RL=-10\log_{10} (0.1489^{2})=16.5421
\end{equation}

\section{Resultados}
\begin{table}[ht]
    \centering
    \begin{tabular}{|l|l|l|l|l|l|l|l|l|}
    \hline
    $Z_L \ [\Omega] $ & $Z_0 \ [\Omega] $ & $\hat{Z}$ & $\hat{Y}$ & $l \ [\lambda] $ & $s \ [\lambda] $ & $\Gamma(0)$ & RL \\ \hline
    47 & 200 & 0.235 & 0.7407 & 0.0722 & 0.411 & -0.6194 & 4.1605 \\ \hline
    200 & 200 & 1.1 & 0.909 & 0.1291 & 0.2347 & 0.0476 & 26.4478  \\ \hline
    200 & 200 & 1.35 & 4.255 & 0.133 & 0.2041 & 0.1489 & 16.5421  \\ \hline
    \end{tabular}
    \caption{Resultados}
    \label{my-label}
\end{table}

\section{Conclusión}
Con la realización de la practica pudimos observar físicamente los fenómenos presentes en el estudio de las líneas de transmisión. Eso se hizo mediante el uso de un kit que permitía realizar mediciones a lo largo de una línea bifilar con ayuda de un generador. Sobre la línea se podía observar la serie de mínimos y máximos presentes en la línea. 
\\ \\
Se comprendido mejor los cálculos realizados al ver la aplicación sobre una línea de transmisión real. Los cálculos se realizaron con ayuda de la carta de Smith, la cual simplifica demasiado las operaciones y ofrece información con una buena aproximación a los cálculos analíticos. 


\end{document}